\documentclass[11pt,a4paper]{article}
\usepackage[utf8]{inputenc}
\usepackage[T1]{fontenc}
\usepackage[spanish,es-minimal,es-noshorthands]{babel}
\usepackage{lmodern}
\linespread{1.22}
\usepackage[a4paper,top=2.4cm,bottom=2.2cm,left=1.5cm,right=1.5cm,headheight=1.1cm,headsep=0.5cm,footskip=1cm]{geometry}
\usepackage{amsmath,amssymb,cancel}
\usepackage{icomma}
\usepackage{graphicx}
\usepackage[table]{xcolor}
\usepackage{tikz}\usetikzlibrary{calc}
\usepackage{pgfplots}\pgfplotsset{compat=1.17}
\usepackage{fontawesome5}
\usepackage{ccicons}
\usepackage{fancyhdr}
\usepackage{enumitem}
\usepackage{booktabs}
\usepackage{tabularx}
\usepackage{needspace}
\usepackage[most]{tcolorbox}
\usepackage{titlesec}
\usepackage[hidelinks]{hyperref}
\hypersetup{pdftitle={Apuntes de teoría. Unidad 8: Distribuciones bidimensionales. 1.º de Bachillerato CCSS},pdfsubject={Matemáticas Aplicadas a las Ciencias Sociales I},pdfkeywords={estadística bidimensional, tablas de doble entrada, nube de puntos, covarianza, correlación, coeficiente de correlación, recta de regresión, estimaciones}}

% colors de la web (FentMates, bloque de Probabilidad y Estadística: morado; es manté el nom fmgreen dels apunts de les altres unitats)
\definecolor{fmnavy}{HTML}{1B2748}
\definecolor{fmgreen}{HTML}{684E88}
\definecolor{fmgreenl}{HTML}{8064A1}
\definecolor{fmyellow}{HTML}{FFC94D}
\definecolor{fmgold}{HTML}{E9A919}
\definecolor{fmaccent}{HTML}{FF7A50}
\definecolor{fmteal}{HTML}{2BB3A3}
\definecolor{fmdef}{HTML}{1D4F73}\definecolor{fmdefbg}{HTML}{EEF4F9}
\definecolor{fmprop}{HTML}{3B4252}\definecolor{fmpropbg}{HTML}{F3F4F7}
\definecolor{fmex}{HTML}{A3203A}\definecolor{fmexbg}{HTML}{FBF0F2}
\definecolor{fmmuted}{HTML}{5D6685}\definecolor{fmball0}{HTML}{E2483D}\definecolor{fmball1}{HTML}{2F6FD6}
\definecolor{fmsurface}{HTML}{EEF1F8}
\definecolor{fmA}{HTML}{3B82F6}\definecolor{fmB}{HTML}{F97316}\definecolor{fmR}{HTML}{E11D48}\definecolor{fmK}{HTML}{16A34A}\definecolor{fmV}{HTML}{7C3AED}

% blocs que no es poden partir: es mesuren i es reserva exactament el seu espai (amb el títol o la capçalera del requadre)
\newsavebox{\fmkb}
\newcommand{\fmmeasure}[2]{\setbox\fmkb=\vbox{\hsize=\dimexpr\linewidth-#1\relax\linewidth=\hsize #2}}
\newcommand{\fmblock}[1]{\fmmeasure{0pt}{#1}\Needspace{\dimexpr\ht\fmkb+\dp\fmkb\relax}#1}
\newcommand{\fmheadblock}[2]{\fmmeasure{0pt}{#2}\Needspace{\dimexpr\ht\fmkb+\dp\fmkb+2.2cm\relax}#1#2}
\newcommand{\fmboxneed}[1]{\fmmeasure{9mm}{#1}\Needspace{\dimexpr\ht\fmkb+\dp\fmkb+1.6cm\relax}}
\newif\iffmstack\newif\iffmfit
\newcommand{\fmfits}[1]{\par\ifdim\dimexpr\pagegoal-\pagetotal\relax<\dimexpr\ht\fmkb+\dp\fmkb+#1\relax\global\fmfitfalse\else\global\fmfittrue\fi}
\newcommand{\fmtwo}[2]{\fmmeasure{0pt}{#1}\fmfits{0pt}\iffmfit #1\else #2\fi}
\newcommand{\fmheadtwo}[3]{\fmmeasure{0pt}{#2}\fmfits{2.2cm}\iffmfit #1#2\else\Needspace{\dimexpr1.6cm+7\baselineskip\relax}#1#3\fi}
\newcommand{\fmtwoin}[2]{\iffmstack #2\else #1\fi}
\newcommand{\fmboxneedalt}[1]{\fmmeasure{9mm}{#1}\fmfits{1.6cm}\iffmfit\global\fmstackfalse\else\global\fmstacktrue\Needspace{8\baselineskip}\fi}
\newcommand{\fmheadboxalt}[2]{\fmmeasure{9mm}{#2}\fmfits{3.4cm}\iffmfit\global\fmstackfalse\else\global\fmstacktrue\Needspace{\dimexpr2.4cm+5\baselineskip\relax}\fi#1}
\newcommand{\fmheadbox}[2]{\fmmeasure{9mm}{#2}\Needspace{\dimexpr\ht\fmkb+\dp\fmkb+3.4cm\relax}#1}
\newcommand{\lt}{<}\newcommand{\gt}{>}
% fórmula en línia que no es parteix mai; si no cap en una línia, es posa centrada
% fórmula en línia; si ocupa més del 80 % de la línia, va centrada (amb la puntuació que la segueix)
\makeatletter
\newcommand{\fmm}[1]{\setbox0=\hbox{$#1$}\ifdim\wd0>0.8\linewidth\def\fmmnext{\fmmdisp{#1}}\else\def\fmmnext{\mbox{\box0}}\fi\fmmnext}
\def\fmmdisp#1{\@ifnextchar.{\fmmdispp{#1}}{\@ifnextchar,{\fmmdispc{#1}}{\@ifnextchar;{\fmmdisps{#1}}{\fmdisp{#1}}}}}
\def\fmmdispp#1.{\fmdisp{#1\,.}}\def\fmmdispc#1,{\fmdisp{#1\,,}}\def\fmmdisps#1;{\fmdisp{#1\,;}}
\makeatother
% fórmula centrada; si és més ampla que la línia, s'encongeix fins que hi cap
\newcommand{\fmdisp}[1]{\[\setbox2=\hbox{$\displaystyle #1$}\ifdim\wd2>\linewidth\resizebox{\linewidth}{!}{\box2}\else\box2\fi\]}
\DeclareUnicodeCharacter{2212}{-}
\DeclareUnicodeCharacter{2225}{\ensuremath{\parallel}}\DeclareUnicodeCharacter{22A5}{\ensuremath{\perp}}\DeclareUnicodeCharacter{2226}{\ensuremath{\nparallel}}
\DeclareUnicodeCharacter{207A}{\ensuremath{^{+}}}\DeclareUnicodeCharacter{207B}{\ensuremath{^{-}}}
\setlength{\parindent}{0pt}
\setlength{\parskip}{12pt}
\setlength{\lineskip}{5pt}\setlength{\lineskiplimit}{5pt}  % més aire al voltant de les fraccions en línia
\setlength{\emergencystretch}{3em}
\setlist{itemsep=4pt,topsep=4pt}
\renewcommand{\thesection}{8.\arabic{section}}
\titleformat{\section}{\color{fmnavy}\Large\bfseries}{\colorbox{fmdef}{\color{white}\strut\ \thesection\ }}{0.6em}{}[{\color{fmdef}\titlerule[1.2pt]}]
\titlespacing*{\section}{0pt}{10pt}{4pt}
\titleformat*{\subsection}{\color{fmgreen}\large\bfseries}
\titlespacing*{\subsection}{0pt}{6pt plus 2pt}{0pt}
\allowdisplaybreaks

% dibuix de la targeta de la unitat a la web (illustrations.js), en blanc
\newcommand{\fmcard}{\begin{tikzpicture}[x=0.062cm,y=-0.062cm,white,line cap=round]
  \begin{scope}[transparency group,opacity=0.15]
  \draw[line width=3pt] (14,104)--(226,104) (24,114)--(24,8);
  \foreach \x/\y in {40/92,58/84,66/90,80/74,92/78,104/64,116/68,128/56,140/60,152/46,166/50,178/38,192/34,206/26}
    \fill (\x,\y) circle (5);
  \draw[line width=4pt,opacity=0.7] (30,98)--(222,22);
  \node[anchor=base west,inner sep=0pt,scale=2.2,font=\large\itshape] at (150,96) {$r\approx0{,}9$};
  \end{scope}
\end{tikzpicture}}
\newcommand{\fmbrand}{\raisebox{-0.25\height}{\includegraphics[height=0.5cm]{logo-fentmates.pdf}}}

% capçalera i peu
\pagestyle{fancy}\fancyhf{}
\fancyhead[L]{\small\fmbrand}
\fancyhead[R]{\small\color{fmmuted}1.º Bachillerato CCSS · Unidad 8: Distribuciones bidimensionales}
\fancyfoot[L]{\scriptsize\color{fmmuted}\href{https://creativecommons.org/licenses/by-nc-sa/4.0/deed.ca}{\ccbyncsa}\ \ FentMates · fentmates.org · CC BY-NC-SA 4.0}
\fancyfoot[R]{\small\color{fmmuted}\thepage}
\renewcommand{\headrulewidth}{0.4pt}\renewcommand{\headrule}{\color{fmmuted!40}\hrule height \headrulewidth}
\fancypagestyle{primera}{\fancyhf{}\renewcommand{\headrulewidth}{0pt}
  \fancyfoot[L]{\scriptsize\color{fmmuted}\href{https://creativecommons.org/licenses/by-nc-sa/4.0/deed.ca}{\ccbyncsa}\ \ FentMates · fentmates.org · CC BY-NC-SA 4.0}
  \fancyfoot[R]{\small\color{fmmuted}\thepage}}

% caixes (mateixos colors que la web)
% quan una caixa es parteix entre dues pàgines, la vora del tall és discontínua
\tcbset{fmbase/.style={enhanced,breakable,arc=3mm,boxrule=1pt,left=3.5mm,right=3.5mm,top=2mm,bottom=2.5mm,before skip=10pt,after skip=10pt,fonttitle=\bfseries,
  frame style={fill=tcbcolframe},
  overlay first={\fill[tcbcolback] ([xshift=1pt]frame.south west) rectangle ([xshift=-1pt,yshift=1.2pt]frame.south east);
    \draw[tcbcolframe,line width=1pt,dash pattern=on 5pt off 3pt] (frame.south west)--(frame.south east);},
  overlay middle={\fill[tcbcolback] ([xshift=1pt]frame.south west) rectangle ([xshift=-1pt,yshift=1.2pt]frame.south east);
    \fill[tcbcolback] ([xshift=1pt]frame.north west) rectangle ([xshift=-1pt,yshift=-1.2pt]frame.north east);
    \draw[tcbcolframe,line width=1pt,dash pattern=on 5pt off 3pt] (frame.south west)--(frame.south east) (frame.north west)--(frame.north east);},
  overlay last={\fill[tcbcolback] ([xshift=1pt]frame.north west) rectangle ([xshift=-1pt,yshift=-1.2pt]frame.north east);
    \draw[tcbcolframe,line width=1pt,dash pattern=on 5pt off 3pt] (frame.north west)--(frame.north east);}}}
\newtcolorbox{fmdef}[1]{fmbase,colback=fmdefbg,colframe=fmdef,coltitle=white,title={#1}}
\newtcolorbox{fmprop}[1]{fmbase,colback=fmpropbg,colframe=fmprop,coltitle=white,title={#1},halign title=center}
\newtcolorbox{fmex}[1]{fmbase,before upper=\raggedright,colback=fmexbg,colframe=fmex,colbacktitle=fmexbg,coltitle=fmex,title={\faPencil*\ #1},
  titlerule=0.4pt,titlerule style={fmex!40,dashed}}
\newtcolorbox{fmwarn}[1]{fmbase,colback=fmsf!6,colframe=fmsf,coltitle=white,title={#1}}
\newtcolorbox{fmnote}[1][]{enhanced,breakable,colback=white,colframe=white,borderline west={3pt}{0pt}{fmaccent},arc=0pt,left=3mm,right=2mm,top=1mm,bottom=1mm,before skip=6pt,after skip=6pt,before upper=\raggedright,#1}
\newcommand{\fmkey}[1]{\colorbox{fmyellow!45}{\strut #1}}
\newcommand{\fmuline}[1]{\textbf{\color{fmaccent}#1}}
\definecolor{fmok}{HTML}{15803D}\definecolor{fmev}{HTML}{2563EB}\definecolor{fmsf}{HTML}{C2410C}\definecolor{fmsi}{HTML}{7C3AED}\definecolor{fmas}{HTML}{A3203A}
\newcommand{\fmbadge}[2]{\tcbox[on line,colback=#1,colframe=#1,arc=2mm,boxsep=0pt,left=4pt,right=4pt,top=1.5pt,bottom=1.5pt]{\color{white}\footnotesize\bfseries #2}}
\newcommand{\fmchip}[1]{\tcbox[on line,colback=fmsf!10,colframe=fmsf!50,arc=2mm,boxsep=1pt,left=4pt,right=4pt]{#1}}
\newcommand{\fmlevel}[2]{\tcbox[on line,colback=white,colframe=fmmuted!40,arc=2mm,boxsep=2pt,halign=center]{\begin{tabular}{c}{\scriptsize\color{fmmuted}\MakeUppercase{#1}}\\#2\end{tabular}}}
% gràfiques
\pgfplotsset{fmaxis/.style={axis lines=middle,grid=major,grid style={fmmuted!15},axis line style={fmmuted,-latex},
  tick label style={font=\scriptsize,color=fmmuted},clip=true,unbounded coords=jump,restrict y to domain=-200:200,
  every axis plot/.append style={line join=round}},
  fmcurve/.style={fmaccent,line width=1.6pt,smooth}}
\tikzset{fmasy/.style={fmmuted,dashed,line width=0.9pt},fmguide/.style={fmaccent,dotted,line width=0.8pt},
  fmdot/.style={circle,fill=fmaccent,inner sep=0pt,minimum size=5pt},fmopen/.style={circle,draw=fmaccent,fill=white,line width=1pt,inner sep=0pt,minimum size=5pt}}
\newcommand{\fmpau}[1]{\tcbox[on line,colback=fmnavy,colframe=fmnavy,arc=2mm,boxsep=0pt,left=4pt,right=4pt,top=1.5pt,bottom=1.5pt]{\color{white}\scriptsize\bfseries\MakeUppercase{#1}}\ }

\begin{document}
\thispagestyle{primera}
\begin{tikzpicture}[remember picture,overlay]
  \fill[fmgreen] (current page.north west) rectangle ([yshift=-10.1cm]current page.north east);
  \begin{scope}
    \clip (current page.north west) rectangle ([yshift=-10.1cm]current page.north east);
    \node[anchor=north west,inner sep=0pt] at ([xshift=6.6cm,yshift=-1.3cm]current page.north west) {\fmcard};
  \end{scope}
  \node[anchor=north west,inner sep=0pt] at ([xshift=1.5cm,yshift=-1.1cm]current page.north west) {\tcbox[colback=white,colframe=white,arc=2.5mm,boxsep=0pt,left=7pt,right=7pt,top=5pt,bottom=5pt]{\includegraphics[height=0.75cm]{logo-fentmates.pdf}}};
\end{tikzpicture}
\vspace*{0.55cm}
\tcbox[colback=white!15!fmgreen,colframe=white!15!fmgreen,arc=2.5mm,boxsep=0pt,left=6pt,right=6pt,top=3pt,bottom=3pt]{\color{white}\small\bfseries Apuntes de teoría · 1.º Bachillerato · Matemáticas Aplicadas a las Ciencias Sociales I}

\vspace{4pt}
{\raggedright\color{white}\fontsize{25}{30}\selectfont\bfseries{\color{fmyellow}Unidad 8} · Distribuciones bidimensionales\par}
\vspace{6pt}
\begin{minipage}{\linewidth}\color{white!88!fmgreen}\small Teoría y ejemplos resueltos paso a paso: distribuciones bidimensionales y tablas de frecuencias (simples y de doble entrada), nube de puntos, correlación y regresión, covarianza, coeficiente de correlación lineal y las dos rectas de regresión, con estimaciones y su fiabilidad.\end{minipage}\par
% l'índex comença sempre a la mateixa altura, per davall de la franja (com en les unitats 5-10)
\vspace*{\dimexpr 8.6cm-\pagetotal\relax}

\begin{tcolorbox}[colback=fmsurface,colframe=fmsurface,arc=3mm,left=4mm,right=4mm,top=2mm,bottom=2mm]
\small\textbf{Índice}\par\smallskip
\begin{minipage}[t]{0.49\linewidth}\small \hyperlink{sec1}{\textbf{\color{fmgreen}8.1}\ \ Distribuciones bidimensionales. Tablas}\\[3pt]
\mbox{\hyperlink{sec2}{\textbf{\color{fmgreen}8.2}\ \ Nube de puntos}}\\[3pt]
\mbox{\hyperlink{sec3}{\textbf{\color{fmgreen}8.3}\ \ Correlación y regresión}}\\[3pt]\end{minipage}\hfill\begin{minipage}[t]{0.49\linewidth}\small \mbox{\hyperlink{sec4}{\textbf{\color{fmgreen}8.4}\ \ Correlación lineal. Covarianza}}\\[3pt]
\mbox{\hyperlink{sec5}{\textbf{\color{fmgreen}8.5}\ \ Coeficiente de correlación}}\\[3pt]
\mbox{\hyperlink{sec6}{\textbf{\color{fmgreen}8.6}\ \ Rectas de regresión}}\\[3pt]\end{minipage}
\end{tcolorbox}

\Needspace{6.2cm}
\section{Distribuciones bidimensionales. Tablas de frecuencias}\hypertarget{sec1}{}

Hasta ahora has estudiado una sola característica de cada individuo (la estadística \textbf{unidimensional}). En esta unidad estudiamos dos a la vez para ver si están relacionadas: si cuando una crece la otra también crece, si decrece o si no hay ninguna relación, y para poder \textbf{estimar} el valor de una a partir de la otra.

\Needspace{5\baselineskip}
\begin{fmdef}{Variable estadística bidimensional}
Cuando en cada individuo estudiamos \textbf{dos} características a la vez, obtenemos una \textbf{variable bidimensional} \fmm{(X,Y)}: cada dato es un par \fmm{(x_i,y_i)}. Por ejemplo, las horas de estudio y la nota de cada alumno, la altura y el peso de cada persona o la temperatura y las ventas de helados de cada día. El objetivo es saber si hay \textbf{relación} entre las dos variables y cómo es de fuerte.
\end{fmdef}

\Needspace{5\baselineskip}
\begin{fmdef}{Dos maneras de organizar los datos}
\begin{itemize}
\item \textbf{Tabla simple} (en filas o columnas): una lista de los \fmm{N} pares. Es la más cómoda cuando hay pocos datos y casi no se repiten.
\item \textbf{Tabla de doble entrada}: en las filas ponemos los valores de \fmm{X}, en las columnas los de \fmm{Y}, y en cada casilla la \textbf{frecuencia conjunta} \fmm{f_{ij}}, el número de individuos con \fmm{X=x_i} e \fmm{Y=y_j}. Las sumas de las filas dan la \textbf{distribución marginal de X} y las de las columnas, la \textbf{distribución marginal de Y}. La suma de todas las casillas es \fmm{N}.
\end{itemize}
\end{fmdef}

\Needspace{5\baselineskip}
\begin{fmex}{Ejemplo: tabla simple}
Una empresa anota el dinero que invierte cada mes en publicidad (\fmm{X}, en miles de euros) y las ventas del mes (\fmm{Y}, en miles de euros):

\begin{center}\renewcommand{\arraystretch}{1.5}\begin{tabular}{ccccccc}\toprule
\textbf{Mes} & \textbf{Ene.} & \textbf{Feb.} & \textbf{Mar.} & \textbf{Abr.} & \textbf{May.} & \textbf{Jun.}\\\midrule
$x_i$ & 1 & 2 & 3 & 4 & 5 & 6\\
$y_i$ & 10 & 14 & 15 & 19 & 21 & 23\\
\bottomrule\end{tabular}\end{center}

Hay \fmm{N=6} pares y ninguno se repite, por eso basta con una tabla simple. Ya se intuye que, cuanto más se invierte en publicidad, más se vende: lo estudiaremos con la nube de puntos (8.2), la covarianza (8.4) y el coeficiente de correlación (8.5), y en el apartado 8.6 calcularemos la recta que permite estimar las ventas.
\end{fmex}

\Needspace{5\baselineskip}
\begin{fmex}{Ejemplo: hermanos y mascotas}
Preguntamos a 20 alumnos el número de hermanos (\fmm{X}) y el número de mascotas que tienen en casa (\fmm{Y}):

\begin{center}\renewcommand{\arraystretch}{1.5}\begin{tabular}{ccccc}\toprule
\textbf{$X\backslash Y$} & \textbf{0} & \textbf{1} & \textbf{2} & \textbf{$f_x$}\\\midrule
0 & 3 & 2 & 0 & 5\\
1 & 2 & 5 & 2 & 9\\
2 & 0 & 2 & 4 & 6\\
$f_y$ & 5 & 9 & 6 & 20\\
\bottomrule\end{tabular}\end{center}

\begin{itemize}
\item La casilla \fmm{(1,1)} vale 5: cinco alumnos tienen un hermano y una mascota.
\item Distribución marginal de \fmm{X}: los valores 0, 1, 2 con frecuencias 5, 9, 6. Su media es \fmm{\bar x=\dfrac{0\cdot5+1\cdot9+2\cdot6}{20}=\dfrac{21}{20}=1{,}05}.
\item Distribución marginal de \fmm{Y}: 0, 1, 2 con frecuencias 5, 9, 6, y también \fmm{\bar y=1{,}05}.
\item Observa que las frecuencias grandes están en la diagonal: los alumnos con más hermanos suelen tener más mascotas.
\end{itemize}
\end{fmex}

\Needspace{5\baselineskip}
\begin{fmprop}{Frecuencias relativas, distribuciones condicionadas y datos agrupados}
\begin{itemize}
\item La \textbf{frecuencia relativa} de una casilla es \fmm{h_{ij}=\dfrac{f_{ij}}{N}} (multiplicada por 100, es un porcentaje). Todas juntas suman 1.
\item Una \textbf{distribución condicionada} es una sola fila (o una sola columna) de la tabla: por ejemplo, cómo se reparte \fmm{Y} entre los individuos que tienen \fmm{X=x_i}. Los porcentajes se calculan respecto del total de esa fila.
\item Si las variables son continuas o tienen muchos valores, se agrupan en \textbf{intervalos} y, para hacer los cálculos, cada intervalo se representa por su \textbf{marca de clase} (el punto medio).
\end{itemize}
\end{fmprop}

\Needspace{5\baselineskip}
\begin{fmex}{Ejemplo: horas de estudio y nota (datos agrupados)}
Las horas de estudio semanales (\fmm{X}) y la nota final (\fmm{Y}) de 40 alumnos, agrupadas en intervalos:

\begin{center}\renewcommand{\arraystretch}{1.5}\begin{tabular}{ccccc}\toprule
\textbf{$X\backslash Y$} & \textbf{$[0,5)$} & \textbf{$[5,7)$} & \textbf{$[7,10]$} & \textbf{$f_x$}\\\midrule
$[0,2)$ & 6 & 3 & 1 & 10\\
$[2,4)$ & 4 & 9 & 5 & 18\\
$[4,6)$ & 1 & 4 & 7 & 12\\
$f_y$ & 11 & 16 & 13 & 40\\
\bottomrule\end{tabular}\end{center}

Marcas de clase: para \fmm{X}, \fmm{1,\ 3,\ 5}; para \fmm{Y}, \fmm{2{,}5,\ 6,\ 8{,}5}.\\[2ex]

Medias marginales: \fmm{\bar x=\dfrac{1\cdot10+3\cdot18+5\cdot12}{40}=\dfrac{124}{40}=3{,}1} horas, \\[2ex]
\fmm{\ \bar y=\dfrac{2{,}5\cdot11+6\cdot16+8{,}5\cdot13}{40}=\dfrac{234}{40}=5{,}85}.\\[2ex]

Frecuencia relativa de la casilla \fmm{([2,4),[5,7))}: \fmm{h=\frac{9}{40}=0{,}225}, es decir, el \fmm{22{,}5\,\%} de los alumnos estudia entre 2 y 4 horas y saca entre un 5 y un 7.

Distribución condicionada a \fmm{X\in[4,6)} (la tercera fila): de los 12 alumnos que estudian más de 4 horas, sacan un 7 o más \fmm{\frac{7}{12}\approx58{,}3\,\%}, mientras que de los que estudian menos de 2 horas solo \fmm{\frac{1}{10}=10\,\%}. Parece que hay relación entre las dos variables.
\end{fmex}

\fmheadboxalt{
\section{Representación gráfica: la nube de puntos}\hypertarget{sec2}{}
}{\par\noindent \begin{minipage}[t]{0.48\linewidth}\vspace{0pt}\raggedright Representamos cada par \fmm{(x_i,y_i)} como un punto del plano. El conjunto de puntos es la \textbf{nube de puntos}. Su forma nos dice, de un vistazo, si hay relación entre las variables:

\begin{itemize}
\item si los puntos se agrupan alrededor de una recta, decimos que hay \textbf{correlación lineal};
\item si la recta sube, la correlación es \textbf{positiva} (directa); si baja, \textbf{negativa} (inversa);
\item cuanto más estrecha es la nube, más \textbf{fuerte} es la relación.
\end{itemize}

El punto \fmm{(\bar x,\bar y)} se llama \textbf{centro de gravedad} de la nube.\end{minipage}\hfill\begin{minipage}[t]{0.48\linewidth}\vspace{0pt}\raggedright \par\medskip\begin{minipage}{\linewidth}{\centering \includegraphics[width=\linewidth,height=4.5cm,keepaspectratio]{figs/a143dbbb6320.pdf}\par}\par\smallskip{\centering\footnotesize Los puntos \fmm{(1,2),(2,3),(3,5),(4,4),(5,6)} y el centro de gravedad \fmm{(3,4)}.\par}\end{minipage}\par\medskip\end{minipage}\par\medskip}
\begin{fmdef}{Nube de puntos (diagrama de dispersión)}
\fmtwoin{\par\noindent \begin{minipage}[t]{0.48\linewidth}\vspace{0pt}\raggedright Representamos cada par \fmm{(x_i,y_i)} como un punto del plano. El conjunto de puntos es la \textbf{nube de puntos}. Su forma nos dice, de un vistazo, si hay relación entre las variables:

\begin{itemize}
\item si los puntos se agrupan alrededor de una recta, decimos que hay \textbf{correlación lineal};
\item si la recta sube, la correlación es \textbf{positiva} (directa); si baja, \textbf{negativa} (inversa);
\item cuanto más estrecha es la nube, más \textbf{fuerte} es la relación.
\end{itemize}

El punto \fmm{(\bar x,\bar y)} se llama \textbf{centro de gravedad} de la nube.\end{minipage}\hfill\begin{minipage}[t]{0.48\linewidth}\vspace{0pt}\raggedright \par\medskip\begin{minipage}{\linewidth}{\centering \includegraphics[width=\linewidth,height=4.5cm,keepaspectratio]{figs/a143dbbb6320.pdf}\par}\par\smallskip{\centering\footnotesize Los puntos \fmm{(1,2),(2,3),(3,5),(4,4),(5,6)} y el centro de gravedad \fmm{(3,4)}.\par}\end{minipage}\par\medskip\end{minipage}\par\medskip}{

Representamos cada par \fmm{(x_i,y_i)} como un punto del plano. El conjunto de puntos es la \textbf{nube de puntos}. Su forma nos dice, de un vistazo, si hay relación entre las variables:

\begin{itemize}
\item si los puntos se agrupan alrededor de una recta, decimos que hay \textbf{correlación lineal};
\item si la recta sube, la correlación es \textbf{positiva} (directa); si baja, \textbf{negativa} (inversa);
\item cuanto más estrecha es la nube, más \textbf{fuerte} es la relación.
\end{itemize}

El punto \fmm{(\bar x,\bar y)} se llama \textbf{centro de gravedad} de la nube.

\begin{center}\begin{minipage}{0.62\linewidth}\centering \par\medskip\begin{minipage}{\linewidth}{\centering \includegraphics[width=\linewidth,height=4.5cm,keepaspectratio]{figs/a143dbbb6320.pdf}\par}\par\smallskip{\centering\footnotesize Los puntos \fmm{(1,2),(2,3),(3,5),(4,4),(5,6)} y el centro de gravedad \fmm{(3,4)}.\par}\end{minipage}\par\medskip\end{minipage}\end{center}
}
\end{fmdef}\global\fmstackfalse

\fmboxneed{\par\noindent \begin{minipage}[t]{0.312\linewidth}\centering {\centering \includegraphics[width=\linewidth,height=3.8cm,keepaspectratio]{figs/0e9e9c81c2e5.pdf}\par}\par\smallskip{\centering\footnotesize Positiva y fuerte · $r\approx0{,}92$\par}\end{minipage}\hfill\begin{minipage}[t]{0.312\linewidth}\centering {\centering \includegraphics[width=\linewidth,height=3.8cm,keepaspectratio]{figs/65f231e9a1ff.pdf}\par}\par\smallskip{\centering\footnotesize Negativa y fuerte · $r\approx-0{,}85$\par}\end{minipage}\hfill\begin{minipage}[t]{0.312\linewidth}\centering {\centering \includegraphics[width=\linewidth,height=3.8cm,keepaspectratio]{figs/9c0877b5d55c.pdf}\par}\par\smallskip{\centering\footnotesize Positiva y débil · $r\approx0{,}45$\par}\end{minipage}\par\medskip}
\begin{fmex}{Cuatro nubes típicas}
\par\noindent \begin{minipage}[t]{0.312\linewidth}\centering {\centering \includegraphics[width=\linewidth,height=3.8cm,keepaspectratio]{figs/0e9e9c81c2e5.pdf}\par}\par\smallskip{\centering\footnotesize Positiva y fuerte · $r\approx0{,}92$\par}\end{minipage}\hfill\begin{minipage}[t]{0.312\linewidth}\centering {\centering \includegraphics[width=\linewidth,height=3.8cm,keepaspectratio]{figs/65f231e9a1ff.pdf}\par}\par\smallskip{\centering\footnotesize Negativa y fuerte · $r\approx-0{,}85$\par}\end{minipage}\hfill\begin{minipage}[t]{0.312\linewidth}\centering {\centering \includegraphics[width=\linewidth,height=3.8cm,keepaspectratio]{figs/9c0877b5d55c.pdf}\par}\par\smallskip{\centering\footnotesize Positiva y débil · $r\approx0{,}45$\par}\end{minipage}\par\medskip
\par\noindent \begin{minipage}[t]{0.312\linewidth}\centering {\centering \includegraphics[width=\linewidth,height=3.8cm,keepaspectratio]{figs/e83c6fcc4dd4.pdf}\par}\par\smallskip{\centering\footnotesize Sin correlación · $r\approx0{,}02$\par}\end{minipage}\hfill\hspace*{0.653\linewidth}\par\medskip

\begin{fmnote}Ejemplos: la altura y el peso de las personas (positiva), el precio de un producto y la cantidad que se vende (negativa), la talla de zapato y la nota de matemáticas (sin relación).\end{fmnote}
\end{fmex}

\Needspace{8\baselineskip}
\section{Correlación y regresión}\hypertarget{sec3}{}

Ante una distribución bidimensional nos hacemos dos preguntas: \textbf{cómo de fuerte} es la relación entre las variables (de eso se encarga la \textbf{correlación}) y \textbf{qué función} la describe, para poder estimar una variable a partir de la otra (de eso se encarga la \textbf{regresión}).

\Needspace{5\baselineskip}
\begin{fmdef}{Tipos de dependencia entre dos variables}
\begin{itemize}
\item \textbf{Dependencia funcional}: el valor de \fmm{X} determina \textbf{exactamente} el de \fmm{Y}; hay una fórmula \fmm{y=f(x)}. Por ejemplo, el precio que pagas por unas manzanas a 2 €/kg y los kilos que compras: \fmm{y=2x}. Todos los puntos están sobre una línea.
\item \textbf{Dependencia estadística} (o aleatoria): hay relación, pero no es exacta. Dos personas de la misma altura no pesan igual, pero en general las más altas pesan más. Los puntos forman una nube alrededor de una línea. Es el caso que estudia esta unidad.
\item \textbf{Independencia}: no hay ninguna relación. Conocer \fmm{X} no nos dice nada de \fmm{Y} (el número de zapato y la nota de matemáticas). La nube no tiene ninguna forma.
\end{itemize}

\par\noindent \begin{minipage}[t]{0.312\linewidth}\centering {\centering \includegraphics[width=\linewidth,height=3.8cm,keepaspectratio]{figs/4403f9bbe3cc.pdf}\par}\par\smallskip{\centering\footnotesize Funcional: \fmm{y=2x}\par}\end{minipage}\hfill\begin{minipage}[t]{0.312\linewidth}\centering {\centering \includegraphics[width=\linewidth,height=3.8cm,keepaspectratio]{figs/eb6f202507aa.pdf}\par}\par\smallskip{\centering\footnotesize Estadística\par}\end{minipage}\hfill\begin{minipage}[t]{0.312\linewidth}\centering {\centering \includegraphics[width=\linewidth,height=3.8cm,keepaspectratio]{figs/c2bef8cb2d3b.pdf}\par}\par\smallskip{\centering\footnotesize Independencia\par}\end{minipage}\par\medskip
\end{fmdef}

\Needspace{5\baselineskip}
\begin{fmprop}{Correlación y regresión}
\begin{itemize}
\item La \textbf{correlación} estudia el \textbf{grado} de dependencia entre las variables. Si la nube se agrupa alrededor de una recta, hablamos de \textbf{correlación lineal}: \textbf{positiva} o directa si la recta sube, \textbf{negativa} o inversa si baja, y más \textbf{fuerte} cuanto más estrecha es la nube. La mediremos con la covarianza (8.4) y, sobre todo, con el coeficiente de correlación (8.5).
\item La \textbf{regresión} busca la función que mejor se ajusta a la nube de puntos para hacer \textbf{estimaciones} o predicciones. Si la función es una recta, es \textbf{regresión lineal} (apartado 8.6). Si la relación tiene forma de parábola, de exponencial\dots{}, es una regresión no lineal (no la estudiaremos).
\end{itemize}

\par\noindent \begin{minipage}[t]{0.312\linewidth}\centering {\centering \includegraphics[width=\linewidth,height=3.8cm,keepaspectratio]{figs/0e9e9c81c2e5.pdf}\par}\par\smallskip{\centering\footnotesize Lineal: la nube sigue una recta\par}\end{minipage}\hfill\begin{minipage}[t]{0.312\linewidth}\centering {\centering \includegraphics[width=\linewidth,height=3.8cm,keepaspectratio]{figs/d04f60868461.pdf}\par}\par\smallskip{\centering\footnotesize Parabólica: el rendimiento de un deportista según la edad\par}\end{minipage}\hfill\begin{minipage}[t]{0.312\linewidth}\centering {\centering \includegraphics[width=\linewidth,height=3.8cm,keepaspectratio]{figs/6e2e2cb093d3.pdf}\par}\par\smallskip{\centering\footnotesize Exponencial: las bacterias de un cultivo a lo largo del tiempo\par}\end{minipage}\par\medskip
\end{fmprop}

\Needspace{5\baselineskip}
\begin{fmprop}{Correlación no significa causalidad}
Que dos variables tengan una correlación fuerte \textbf{no} significa que una sea la causa de la otra. Por ejemplo, en los pueblos de la costa, las ventas de helados y el número de quemaduras solares tienen una correlación positiva fuerte, pero los helados no provocan quemaduras: las dos dependen de una \textbf{tercera variable}, el calor (y el número de turistas). También hay correlaciones por pura casualidad.
\end{fmprop}

\Needspace{5.7cm}
\section{Correlación lineal. Covarianza}\hypertarget{sec4}{}
\begin{fmdef}{Parámetros marginales}
Cada variable por separado tiene su media, su varianza y su desviación típica (son los parámetros de las distribuciones marginales):

\fmdisp{\bar x=\frac{\sum x_i}{N}\qquad\bar y=\frac{\sum y_i}{N}\qquad\sigma_x^2=\frac{\sum x_i^2}{N}-\bar x^2\qquad\sigma_y^2=\frac{\sum y_i^2}{N}-\bar y^2}

Si los datos están en una tabla de doble entrada, cada producto se multiplica por la frecuencia: \fmm{\bar x=\frac{\sum x_if_i}{N}}, etc.
\end{fmdef}

\Needspace{5\baselineskip}
\begin{fmprop}{Covarianza}
La \textbf{covarianza} mide cómo varían las dos variables \textbf{conjuntamente}:

\fmdisp{\sigma_{xy}=\frac{\sum(x_i-\bar x)(y_i-\bar y)}{N}=\frac{\sum x_i\,y_i}{N}-\bar x\,\bar y\qquad\left(\text{con frecuencias: }\ \frac{\sum x_i\,y_j\,f_{ij}}{N}-\bar x\,\bar y\right)}

\begin{itemize}
\item \fmm{\sigma_{xy}\gt0}: cuando una variable crece, la otra también tiende a crecer. La mayoría de puntos están en los cuadrantes 1º y 3º respecto de \fmm{(\bar x,\bar y)}, donde \fmm{(x_i-\bar x)(y_i-\bar y)\gt0}.
\item \fmm{\sigma_{xy}\lt0}: cuando una crece, la otra tiende a decrecer (cuadrantes 2º y 4º).
\item \fmm{\sigma_{xy}\approx0}: no hay relación lineal.
\end{itemize}

\begin{fmnote}La covarianza depende de las unidades (si pasamos de metros a centímetros, se multiplica por 100), por eso no nos dice lo \textbf{fuerte} que es la relación. Para eso usaremos el coeficiente de correlación (apartado 8.5).\end{fmnote}
\end{fmprop}

\Needspace{5\baselineskip}
\begin{fmex}{Ejemplos}
\begin{enumerate}[label=\textbf{\arabic*.},leftmargin=*]
\item Los puntos del apartado 8.2: \fmm{(1,2),(2,3),(3,5),(4,4),(5,6)}. Hacemos una tabla con las columnas que necesitamos:

\begin{center}\renewcommand{\arraystretch}{1.5}\begin{tabular}{ccccc}\toprule
\textbf{$x_i$} & \textbf{$y_i$} & \textbf{$x_i^2$} & \textbf{$y_i^2$} & \textbf{$x_iy_i$}\\\midrule
1 & 2 & 1 & 4 & 2\\
2 & 3 & 4 & 9 & 6\\
3 & 5 & 9 & 25 & 15\\
4 & 4 & 16 & 16 & 16\\
5 & 6 & 25 & 36 & 30\\
15 & 20 & 55 & 90 & 69\\
\bottomrule\end{tabular}\end{center}

\fmm{\bar x=\dfrac{15}{5}=3\qquad\bar y=\dfrac{20}{5}=4}

\fmm{\sigma_x^2=\dfrac{55}{5}-3^2=2\qquad\sigma_y^2=\dfrac{90}{5}-4^2=2\qquad\sigma_x=\sigma_y=\sqrt2\approx1{,}41}

\fmm{\sigma_{xy}=\dfrac{69}{5}-3\cdot4=13{,}8-12=1{,}8\gt0}: relación positiva.
\item Con la tabla de doble entrada del apartado 8.1 (hermanos y mascotas, \fmm{\bar x=\bar y=1{,}05}), sumamos \fmm{x_i\,y_j\,f_{ij}} en las casillas (las que tienen un 0 no suman):

\fmm{\sum x_iy_jf_{ij}=1\cdot1\cdot5+1\cdot2\cdot2+2\cdot1\cdot2+2\cdot2\cdot4=5+4+4+16=29}

\fmm{\sigma_{xy}=\dfrac{29}{20}-1{,}05\cdot1{,}05=1{,}45-1{,}1025=0{,}3475\gt0}

\fmm{\sum x_i^2f_x=0\cdot5+1\cdot9+4\cdot6=33\Rightarrow\sigma_x^2=\dfrac{33}{20}-1{,}05^2=0{,}5475}. Igualmente, \fmm{\sigma_y^2=0{,}5475}.
\item Publicidad (\fmm{X}) y ventas (\fmm{Y}) del apartado 8.1, en miles de euros:

\begin{center}\renewcommand{\arraystretch}{1.5}\begin{tabular}{ccccc}\toprule
\textbf{$x_i$} & \textbf{$y_i$} & \textbf{$x_i^2$} & \textbf{$y_i^2$} & \textbf{$x_iy_i$}\\\midrule
1 & 10 & 1 & 100 & 10\\
2 & 14 & 4 & 196 & 28\\
3 & 15 & 9 & 225 & 45\\
4 & 19 & 16 & 361 & 76\\
5 & 21 & 25 & 441 & 105\\
6 & 23 & 36 & 529 & 138\\
21 & 102 & 91 & 1852 & 402\\
\bottomrule\end{tabular}\end{center}

\fmm{\bar x=\dfrac{21}{6}=3{,}5\qquad\bar y=\dfrac{102}{6}=17}

\fmm{\sigma_x^2=\dfrac{91}{6}-3{,}5^2\approx2{,}917\qquad\sigma_y^2=\dfrac{1852}{6}-17^2\approx19{,}667}

\fmm{\sigma_{xy}=\dfrac{402}{6}-3{,}5\cdot17=67-59{,}5=7{,}5\gt0}: cuando se invierte más en publicidad, las ventas tienden a ser más altas.
\end{enumerate}
\end{fmex}

\Needspace{5.7cm}
\section{Coeficiente de correlación lineal}\hypertarget{sec5}{}
\begin{fmdef}{Coeficiente de correlación lineal de Pearson}

\fmdisp{r=\frac{\sigma_{xy}}{\sigma_x\,\sigma_y}}

Como dividimos por las desviaciones típicas, \fmm{r} no tiene unidades y siempre cumple \fmm{-1\leq r\leq1}. Tiene el mismo signo que la covarianza.
\end{fmdef}

\Needspace{5\baselineskip}
\begin{fmprop}{Cómo se interpreta}
\begin{itemize}
\item \textbf{El signo} indica el sentido: \fmm{r\gt0}, correlación positiva (directa); \fmm{r\lt0}, negativa (inversa).
\item \textbf{El valor absoluto} indica la fuerza: \fmm{|r|=1}, los puntos están exactamente sobre una recta (dependencia \textbf{funcional}); \fmm{|r|} cerca de 1 (por ejemplo, más de 0,8), correlación \textbf{fuerte}; entre 0,5 y 0,8, \textbf{moderada}; por debajo de 0,5, \textbf{débil}; \fmm{r\approx0}, no hay correlación lineal.
\end{itemize}

\begin{itemize}
\item \fmm{r} \textbf{no depende de las unidades}: si medimos en centímetros en lugar de metros, \fmm{r} no cambia (la covarianza, sí).
\item Si \fmm{r=0} decimos que las variables están \textbf{incorreladas}: no hay relación \textbf{lineal}, pero puede haberla de otro tipo (véase más abajo).
\end{itemize}

\begin{fmnote}Estos límites (0,8 y 0,5) son orientativos: no hay una frontera exacta entre una correlación fuerte y una moderada.\end{fmnote}
\end{fmprop}

\Needspace{5\baselineskip}
\begin{fmex}{Ejemplos}
\begin{enumerate}[label=\textbf{\arabic*.},leftmargin=*]
\item Los puntos \fmm{(1,2),(2,3),(3,5),(4,4),(5,6)} del apartado 8.4: \fmm{\sigma_{xy}=1{,}8} y \fmm{\sigma_x=\sigma_y=\sqrt2}.

\fmdisp{r=\frac{1{,}8}{\sqrt2\cdot\sqrt2}=\frac{1{,}8}{2}=0{,}9}

          Correlación lineal positiva y fuerte.
\item Hermanos y mascotas (apartado 8.4): \fmm{\ r=\dfrac{0{,}3475}{\sqrt{0{,}5475}\cdot\sqrt{0{,}5475}}=\dfrac{0{,}3475}{0{,}5475}\approx0{,}63}: correlación positiva moderada.
\item El precio (€) y las unidades vendidas de un producto en 5 tiendas: \fmm{(2,40),(3,34),(4,30),(5,25),(6,21)}. Ahora \fmm{\sum x_i=20}, \fmm{\sum y_i=150}, \fmm{\sum x_i^2=90}, \fmm{\sum y_i^2=4\,722} y \fmm{\sum x_iy_i=553}.

\fmm{\bar x=4,\ \bar y=30,\ \sigma_x^2=\dfrac{90}{5}-16=2,\ \sigma_y^2=\dfrac{4\,722}{5}-900=44{,}4}

\fmm{\sigma_{xy}=\dfrac{553}{5}-4\cdot30=110{,}6-120=-9{,}4}

\fmm{r=\dfrac{-9{,}4}{\sqrt2\cdot\sqrt{44{,}4}}\approx\dfrac{-9{,}4}{9{,}423}\approx-0{,}998}: correlación negativa muy fuerte (cuanto más caro, menos se vende).
\item Publicidad y ventas (apartado 8.4): \fmm{\sigma_{xy}=7{,}5}, \fmm{\sigma_x^2\approx2{,}917}, \fmm{\sigma_y^2\approx19{,}667}.

\fmdisp{r=\frac{7{,}5}{\sqrt{2{,}917}\cdot\sqrt{19{,}667}}\approx\frac{7{,}5}{1{,}708\cdot4{,}435}\approx0{,}99}

          Correlación positiva muy fuerte: tiene sentido usar una recta para estimar las ventas (apartado 8.6).
\end{enumerate}
\end{fmex}

\fmboxneed{\par\noindent \begin{minipage}[t]{0.312\linewidth}\centering {\centering \includegraphics[width=\linewidth,height=3.8cm,keepaspectratio]{figs/d04f60868461.pdf}\par}\par\smallskip{\centering\footnotesize Parabólica: el rendimiento de un deportista según la edad. Crece y después decrece. $r\approx0$\par}\end{minipage}\hfill\begin{minipage}[t]{0.312\linewidth}\centering {\centering \includegraphics[width=\linewidth,height=3.8cm,keepaspectratio]{figs/6e2e2cb093d3.pdf}\par}\par\smallskip{\centering\footnotesize Exponencial: el número de bacterias de un cultivo a lo largo del tiempo. $r\approx0{,}89$\par}\end{minipage}\hfill\begin{minipage}[t]{0.312\linewidth}\centering {\centering \includegraphics[width=\linewidth,height=3.8cm,keepaspectratio]{figs/f507e832ac94.pdf}\par}\par\smallskip{\centering\footnotesize De saturación: lo que recuerdas de un tema según las horas que le dedicas. $r\approx0{,}85$\par}\end{minipage}\par\medskip}
\begin{fmex}{Cuidado: $r$ solo mide la relación lineal}
\par\noindent \begin{minipage}[t]{0.312\linewidth}\centering {\centering \includegraphics[width=\linewidth,height=3.8cm,keepaspectratio]{figs/d04f60868461.pdf}\par}\par\smallskip{\centering\footnotesize Parabólica: el rendimiento de un deportista según la edad. Crece y después decrece. $r\approx0$\par}\end{minipage}\hfill\begin{minipage}[t]{0.312\linewidth}\centering {\centering \includegraphics[width=\linewidth,height=3.8cm,keepaspectratio]{figs/6e2e2cb093d3.pdf}\par}\par\smallskip{\centering\footnotesize Exponencial: el número de bacterias de un cultivo a lo largo del tiempo. $r\approx0{,}89$\par}\end{minipage}\hfill\begin{minipage}[t]{0.312\linewidth}\centering {\centering \includegraphics[width=\linewidth,height=3.8cm,keepaspectratio]{figs/f507e832ac94.pdf}\par}\par\smallskip{\centering\footnotesize De saturación: lo que recuerdas de un tema según las horas que le dedicas. $r\approx0{,}85$\par}\end{minipage}\par\medskip
\par\noindent \begin{minipage}[t]{0.312\linewidth}\centering {\centering \includegraphics[width=\linewidth,height=3.8cm,keepaspectratio]{figs/54977b7b6c4f.pdf}\par}\par\smallskip{\centering\footnotesize Los puntos forman una circunferencia: hay una relación perfecta, pero \fmm{r=0}. $r\approx0$\par}\end{minipage}\hfill\hspace*{0.653\linewidth}\par\medskip

\begin{fmnote}Fíjate en la parábola y en la circunferencia: \fmm{r} es prácticamente 0 y, en cambio, \fmm{Y} depende claramente de \fmm{X}. Por eso \textbf{siempre} hay que mirar la nube de puntos antes de interpretar \fmm{r}. En la exponencial y en la de saturación, \fmm{r} es alto, pero los puntos no siguen una recta.\end{fmnote}
\end{fmex}

\fmboxneedalt{\par\noindent \begin{minipage}[t]{0.48\linewidth}\vspace{0pt}\raggedright Un solo punto muy alejado de los demás (en rojo) puede cambiar mucho el coeficiente. Aquí los siete puntos de abajo no tienen ninguna relación entre ellos, pero el punto atípico hace que \fmm{r} salga muy alto. Antes de calcular, hay que revisar si hay datos erróneos o excepcionales.\end{minipage}\hfill\begin{minipage}[t]{0.48\linewidth}\vspace{0pt}\raggedright \par\medskip\begin{minipage}{\linewidth}{\centering \includegraphics[width=\linewidth,height=4.5cm,keepaspectratio]{figs/86bff594a541.pdf}\par}\par\smallskip{\centering\footnotesize $r\approx0{,}98$\par}\end{minipage}\par\medskip\end{minipage}\par\medskip}
\begin{fmex}{Valores atípicos}
\fmtwoin{\par\noindent \begin{minipage}[t]{0.48\linewidth}\vspace{0pt}\raggedright Un solo punto muy alejado de los demás (en rojo) puede cambiar mucho el coeficiente. Aquí los siete puntos de abajo no tienen ninguna relación entre ellos, pero el punto atípico hace que \fmm{r} salga muy alto. Antes de calcular, hay que revisar si hay datos erróneos o excepcionales.\end{minipage}\hfill\begin{minipage}[t]{0.48\linewidth}\vspace{0pt}\raggedright \par\medskip\begin{minipage}{\linewidth}{\centering \includegraphics[width=\linewidth,height=4.5cm,keepaspectratio]{figs/86bff594a541.pdf}\par}\par\smallskip{\centering\footnotesize $r\approx0{,}98$\par}\end{minipage}\par\medskip\end{minipage}\par\medskip}{

Un solo punto muy alejado de los demás (en rojo) puede cambiar mucho el coeficiente. Aquí los siete puntos de abajo no tienen ninguna relación entre ellos, pero el punto atípico hace que \fmm{r} salga muy alto. Antes de calcular, hay que revisar si hay datos erróneos o excepcionales.

\begin{center}\begin{minipage}{0.62\linewidth}\centering \par\medskip\begin{minipage}{\linewidth}{\centering \includegraphics[width=\linewidth,height=4.5cm,keepaspectratio]{figs/86bff594a541.pdf}\par}\par\smallskip{\centering\footnotesize $r\approx0{,}98$\par}\end{minipage}\par\medskip\end{minipage}\end{center}
}
\end{fmex}\global\fmstackfalse

\Needspace{8\baselineskip}
\section{Rectas de regresión}\hypertarget{sec6}{}

Cuando hay correlación lineal, podemos dibujar la recta que mejor se ajusta a la nube de puntos y usarla para \textbf{estimar} valores de una variable a partir de la otra. Como no es lo mismo estimar \fmm{y} conocida \fmm{x} que estimar \fmm{x} conocida \fmm{y}, hay \textbf{dos} rectas de regresión.

\Needspace{5\baselineskip}
\begin{fmdef}{Recta de regresión de Y sobre X}
Es la recta que hace mínima la suma de los cuadrados de las distancias \textbf{verticales} de los puntos a la recta (método de los \textbf{mínimos cuadrados}). Pasa por el centro de gravedad \fmm{(\bar x,\bar y)} y su pendiente es el \textbf{coeficiente de regresión de Y sobre X}, \fmm{m_{yx}=\frac{\sigma_{xy}}{\sigma_x^2}}:

\fmdisp{y-\bar y=\frac{\sigma_{xy}}{\sigma_x^2}\,(x-\bar x)}

Sirve para \textbf{estimar \fmm{y}} cuando conocemos \fmm{x}. El coeficiente de regresión dice cuánto aumenta (o disminuye, si es negativo) \fmm{y}, de media, por cada unidad que aumenta \fmm{x}.
\end{fmdef}

\Needspace{5\baselineskip}
\begin{fmdef}{Recta de regresión de X sobre Y}
Hace mínima la suma de los cuadrados de las distancias \textbf{horizontales}. También pasa por \fmm{(\bar x,\bar y)} y su coeficiente es \fmm{m_{xy}=\frac{\sigma_{xy}}{\sigma_y^2}}:

\fmdisp{x-\bar x=\frac{\sigma_{xy}}{\sigma_y^2}\,(y-\bar y)}

Sirve para \textbf{estimar \fmm{x}} cuando conocemos \fmm{y}. Cuidado: para estimar \fmm{x} \textbf{no} se despeja \fmm{x} de la recta de Y sobre X; hay que usar esta otra recta.
\end{fmdef}

\Needspace{5\baselineskip}
\begin{fmprop}{Propiedades de las dos rectas}
\begin{itemize}
\item Las dos se cortan en el centro de gravedad \fmm{(\bar x,\bar y)}.
\item Los dos coeficientes de regresión tienen el signo de \fmm{\sigma_{xy}} (y de \fmm{r}): las dos rectas son crecientes o las dos decrecientes.
\item El producto de los coeficientes es el cuadrado del coeficiente de correlación: \fmm{\ m_{yx}\cdot m_{xy}=\dfrac{\sigma_{xy}}{\sigma_x^2}\cdot\dfrac{\sigma_{xy}}{\sigma_y^2}=\dfrac{\sigma_{xy}^2}{\sigma_x^2\sigma_y^2}=r^2}. Por eso \fmm{r=\pm\sqrt{m_{yx}\,m_{xy}}}, con el signo de los coeficientes.
\item Si \fmm{|r|=1}, las dos rectas \textbf{coinciden} (todos los puntos están sobre la recta). Cuanto más cerca de 0 está \fmm{r}, más se abre el ángulo entre ellas; si \fmm{r=0}, son \fmm{y=\bar y} y \fmm{x=\bar x}: \textbf{perpendiculares}.
\item \textbf{¿Cuál es cuál?} Si escribimos las dos en la forma \fmm{y=ax+b}, la de Y sobre X es la de pendiente menor en valor absoluto (la más «plana»): su pendiente es \fmm{m_{yx}} y la de la otra es \fmm{\frac{1}{m_{xy}}}, y el cociente \fmm{m_{yx}:\frac{1}{m_{xy}}=r^2} no puede pasar de 1.
\end{itemize}
\end{fmprop}

\Needspace{5\baselineskip}
\begin{fmprop}{Fiabilidad de las estimaciones}
\begin{itemize}
\item Cuanto más cerca de 1 esté \fmm{|r|}, más fiable es la estimación. Si \fmm{|r|} es pequeño, las rectas no sirven para estimar.
\item Solo son fiables las estimaciones para valores \textbf{dentro} (o cerca) del intervalo de los datos. Fuera de él (\textbf{extrapolar}) la relación puede cambiar.
\end{itemize}
\end{fmprop}

\Needspace{5\baselineskip}
\begin{fmex}{Ejemplos}
\begin{enumerate}[label=\textbf{\arabic*.},leftmargin=*]
\item Los puntos \fmm{(1,2),(2,3),(3,5),(4,4),(5,6)} de los apartados 8.4 y 8.5: \fmm{\bar x=3}, \fmm{\bar y=4}, \fmm{\sigma_x^2=\sigma_y^2=2}, \fmm{\sigma_{xy}=1{,}8}, \fmm{r=0{,}9}.

Recta de Y sobre X: \fmm{\ y-4=\dfrac{1{,}8}{2}(x-3)\Rightarrow y=0{,}9x+1{,}3}

Estimación para \fmm{x=4{,}5}: \fmm{\ \hat y=0{,}9\cdot4{,}5+1{,}3=5{,}35}. Es fiable: \fmm{r=0{,}9} y \fmm{4{,}5} está dentro del intervalo \fmm{[1,5]}.

Recta de X sobre Y: \fmm{\ x-3=\dfrac{1{,}8}{2}(y-4)\Rightarrow x=0{,}9y-0{,}6}

\par\medskip\begin{minipage}{\linewidth}{\centering \includegraphics[width=7cm,height=5.2cm,keepaspectratio]{figs/0ec023f02656.pdf}\par}\end{minipage}\par\medskip
\item Publicidad (\fmm{X}) y ventas (\fmm{Y}), en miles de euros (apartados 8.4 y 8.5): \fmm{\bar x=3{,}5}, \fmm{\bar y=17}, \fmm{\sigma_x^2=\frac{35}{12}\approx2{,}917}, \fmm{\sigma_y^2=\frac{59}{3}\approx19{,}667}, \fmm{\sigma_{xy}=7{,}5}, \fmm{r\approx0{,}99}.

Recta de Y sobre X: \fmm{\ m_{yx}=\dfrac{7{,}5}{35/12}=\dfrac{18}{7}\approx2{,}571}, \fmm{\ y-17=2{,}571\,(x-3{,}5)\Rightarrow y\approx2{,}571x+8}. Por cada 1000 € más de publicidad, las ventas suben, de media, unos 2571 €.

Si se invierten 4500 € (\fmm{x=4{,}5}), se estiman unas ventas de \fmm{\ \hat y=2{,}571\cdot4{,}5+8\approx19{,}57}, es decir, unos 19 570 €.

Recta de X sobre Y: \fmm{\ m_{xy}=\dfrac{7{,}5}{59/3}=\dfrac{45}{118}\approx0{,}381}, \fmm{\ x-3{,}5=0{,}381\,(y-17)\Rightarrow x\approx0{,}381y-2{,}983}.

Para vender 20 000 € (\fmm{y=20}) habría que invertir \fmm{\ \hat x=0{,}381\cdot20-2{,}983\approx4{,}64}, unos 4640 € en publicidad. Las dos estimaciones son fiables: \fmm{r\approx0{,}99} y los valores están dentro de los datos.

Comprobación: \fmm{\ m_{yx}\cdot m_{xy}=\dfrac{18}{7}\cdot\dfrac{45}{118}\approx0{,}981=r^2\Rightarrow r\approx0{,}99}.

\par\medskip\begin{minipage}{\linewidth}{\centering \includegraphics[width=7cm,height=5.2cm,keepaspectratio]{figs/0e571202be8f.pdf}\par}\par\smallskip{\centering\footnotesize En azul, la recta de Y sobre X; en verde, la de X sobre Y. Se cortan en \fmm{(3{,}5;\ 17)}.\par}\end{minipage}\par\medskip
\item Precio (€) y unidades vendidas (apartado 8.5): \fmm{\bar x=4}, \fmm{\bar y=30}, \fmm{\sigma_x^2=2}, \fmm{\sigma_{xy}=-9{,}4}, \fmm{r\approx-0{,}998}.

\fmm{y-30=\dfrac{-9{,}4}{2}(x-4)\Rightarrow y=-4{,}7x+48{,}8}

Si el precio es 4,50 €, se estima que se venderán \fmm{\hat y=-4{,}7\cdot4{,}5+48{,}8=27{,}65\approx28} unidades (muy fiable). En cambio, para un precio de 12 € saldría \fmm{-7{,}6} unidades: no tiene sentido, porque 12 está muy lejos de los datos.
\item Las rectas de regresión de una distribución son \fmm{x+2y=10} y \fmm{2x+y=11}. Halla el centro de gravedad, di cuál es cada recta y calcula \fmm{r}.

Centro de gravedad: resolvemos el sistema formado por las dos rectas. \fmm{\ \left\{\begin{array}{l}x+2y=10\\2x+y=11\end{array}\right.\Rightarrow x=4,\ y=3}. Por tanto, \fmm{(\bar x,\bar y)=(4,3)}.

En la forma \fmm{y=ax+b}: \fmm{y=-0{,}5x+5} e \fmm{y=-2x+11}. La más plana es la primera: es la de \textbf{Y sobre X}, con \fmm{m_{yx}=-0{,}5}. La segunda, despejando \fmm{x}: \fmm{x=-0{,}5y+5{,}5}, es la de \textbf{X sobre Y}, con \fmm{m_{xy}=-0{,}5}.

\fmm{r^2=m_{yx}\cdot m_{xy}=(-0{,}5)(-0{,}5)=0{,}25\Rightarrow r=-0{,}5} (negativo, como los coeficientes): correlación negativa moderada.

Si lo hubiéramos hecho al revés, saldría \fmm{r^2=(-2)(-2)=4\gt1}: imposible. Así se comprueba que la elección es la buena.
\item De una distribución sabemos que \fmm{\bar x=10}, \fmm{\bar y=20}, \fmm{\sigma_x=2}, \fmm{\sigma_y=5} y \fmm{r=0{,}8}. Halla las dos rectas de regresión.

Primero, la covarianza: \fmm{\ \sigma_{xy}=r\,\sigma_x\,\sigma_y=0{,}8\cdot2\cdot5=8}.

Y sobre X: \fmm{\ y-20=\dfrac{8}{4}(x-10)\Rightarrow y=2x}. \fmm{\quad} X sobre Y: \fmm{\ x-10=\dfrac{8}{25}(y-20)\Rightarrow x=0{,}32y+3{,}6}.

Comprobación: \fmm{\ 2\cdot0{,}32=0{,}64=0{,}8^2}.
\end{enumerate}
\end{fmex}

  
\end{document}
